\documentclass[11pt]{article}
\usepackage[a4paper,margin=28mm]{geometry}
\usepackage{amsmath,amssymb,amsthm}
\usepackage[colorlinks=true,linkcolor=blue,urlcolor=blue,citecolor=blue]{hyperref}
\newtheorem{theorem}{Theorem}
\newcommand{\Q}{\mathbb Q}
\newcommand{\Z}{\mathbb Z}
\newcommand{\F}{\mathbb F}
\DeclareMathOperator{\End}{End}
\DeclareMathOperator{\Gal}{Gal}
\DeclareMathOperator{\GL}{GL}
\DeclareMathOperator{\PGL}{PGL}
\title{Prime torsion on geometrically split abelian surfaces}
\author{}
\date{20 October 2026}

\begin{document}
\maketitle

\begin{theorem}\label{thm:main}
Let \(A/\Q\) be a geometrically split abelian surface. Then
\[
 A(\Q)[23]=A(\Q)[29]=0.
\]
In particular, neither \([23]\) nor \([29]\) occurs as the rational
torsion subgroup of a geometrically split genus-\(2\) Jacobian over
\(\Q\).
\end{theorem}

Here ``geometrically split'' means that \(A_{\overline{\Q}}\) is
isogenous to a product of elliptic curves. We impose no restriction on
the degree of this isogeny. The theorem does not require a principal
polarization.

\section{Inputs on elliptic torsion and potentially good reduction}

For a positive integer \(d\), let \(S(d)\) be the set of primes that
occur as the order of a torsion point on an elliptic curve over a number
field of degree \(d\). We use the following results.
\begin{enumerate}
\item Derickx--Kamienny--Stein--Stoll \cite[Theorem 1.2]{DKSS} prove
that
\[
 S(4)=\{2,3,5,7,11,13,17\}.
\]
Consequently an elliptic curve over a number field of degree \(1\),
\(2\), or \(4\) has no point of order \(23\) or \(29\). For degrees
\(1\) and \(2\), extend the field to one of degree \(4\).
\item Khawaja \cite[Theorem 1.2]{Khawaja} proves that
\[
 S(8)=\{2,3,5,7,11,13,17,19,23\}.
\]
In particular, an elliptic curve over a degree-\(8\) number field has
no point of order \(29\).
\item Clark \cite[Theorem 5]{Clark} proves that, if an abelian surface
over \(\Q\) has everywhere potentially good reduction, the order of
its rational torsion subgroup belongs to
\[
 \{1,\ldots,16,18,19,20,22,24,25,28,30,36,48,60,72\}.
\]
None of these integers is divisible by \(23\) or \(29\).
\end{enumerate}
We also use the standard properties of Tate modules and geometric
endomorphisms, and the independence of the semisimplified residual
representation from the choice of a stable lattice.

\section{Nonisogenous factors and the CM case}

Fix \(p\in\{23,29\}\), and suppose that \(P\in A(\Q)\) has order \(p\).
First suppose that
\(A_{\overline{\Q}}\) is isogenous to \(E_1\times E_2\), where \(E_1\)
and \(E_2\) are not geometrically isogenous. There are exactly two
geometric elliptic subvarieties of \(A\). Galois permutes them, so
both are defined over a field \(K\) with \([K:\Q]\leq 2\).
Choose one of them, \(E\), and set \(F=A_K/E\). There is an exact
sequence over \(K\)
\[
 0\longrightarrow E\longrightarrow A_K\longrightarrow F
 \longrightarrow 0.
\]
The image of \(P\) in \(F(K)\) either has order \(p\) or is zero. In
the latter case \(P\in E(K)\). Both possibilities contradict the
degree-\(4\) elliptic torsion bound above.

Next suppose that \(A_{\overline{\Q}}\) is isogenous to \(E^2\), where
\(E\) has complex multiplication. CM elliptic curves have everywhere
potentially good reduction. This property is preserved by products
and isogenies, and descends from a finite extension of the ground
field. Hence \(A\) has everywhere potentially good reduction, and
Clark's theorem gives a contradiction.

It remains to consider the case where
\(A_{\overline{\Q}}\) is isogenous to \(E^2\) for a non-CM elliptic
curve \(E\). Then
\[
 D=\End^0(A_{\overline{\Q}})\simeq M_2(\Q).
\]

\section{The field of definition of the endomorphisms}

Let \(L/\Q\) be the smallest field over which all geometric
endomorphisms of \(A\) are defined. It is a finite Galois extension.
Put \(\Gamma=\Gal(L/\Q)\). By Skolem--Noether, its faithful action
on \(D\) identifies \(\Gamma\) with a finite subgroup of
\[
 \operatorname{Aut}_{\Q}(M_2(\Q))=\PGL_2(\Q).
\]

A finite subgroup of \(\PGL_2(\mathbb R)\) is cyclic or dihedral:
its action on the hyperbolic plane has a fixed point, and the action
on the tangent plane embeds it in \(O(2)\). Let an element of
\(\PGL_2(\Q)\) have finite order \(n>2\), and let \(M\) be a
representative. A non-semisimple \(M\) is a scalar times a nontrivial
unipotent matrix and has infinite order in \(\PGL_2(\Q)\). Hence \(M\)
is semisimple, and the ratio of its eigenvalues is a primitive
\(n\)-th root of unity \(\zeta\).
Thus
\[
 \frac{(\operatorname{tr}M)^2}{\det M}
 =2+\zeta+\zeta^{-1}\in\Q,
\]
which implies \(n\in\{3,4,6\}\). It follows that \(\Gamma\) is cyclic
of order \(1,2,3,4,\) or \(6\), or dihedral of order \(4,6,8,\) or
\(12\). Here the dihedral group of order \(4\) is the Klein four
group. In particular,
\[
 [L:\Q]\in\{1,2,3,4,6,8,12\}.
\]

Every geometric elliptic subvariety of \(A\) is defined over \(L\).
Indeed, it is the image of an integral multiple of an idempotent in
\(D\), and that endomorphism is defined over \(L\). If
\([L:\Q]\in\{1,2,4\}\), the elliptic exact sequence used in the
previous section, now over \(L\), excludes \(P\). This includes the
Klein four case. If \(p=29\), the degree-\(8\) elliptic torsion bound
also excludes the case \([L:\Q]=8\).

In every remaining case, \(\Gamma\) has a normal cyclic subgroup
\(\langle r\rangle\) of order \(n\), with
\[
 \begin{array}{c|c}
 p & n\\ \hline
 23 & 3,4,\text{ or }6\\
 29 & 3\text{ or }6 .
 \end{array}
\]
For a cyclic group take the whole group; for a dihedral group take
the rotation subgroup.

\section{A stable maximal order}

Set \(k=\Q_p\) and \(\mathcal O=\Z_p\), and identify
\(D\otimes_{\Q}k\) with \(M_2(k)\).
Choose a matrix representative of \(r\), and divide it by its
nonzero trace. The resulting matrix \(R\) has characteristic
polynomial
\[
 X^2-X+\frac1c,\qquad
 c=
 \begin{cases}
 1,&n=3,\\
 2,&n=4,\\
 3,&n=6.
 \end{cases}
\]
Its discriminant is respectively \(-3,-1,\) or \(-1/3\).
Each is a nonsquare unit modulo \(p\) in the cases retained above:
\(23\equiv-1\pmod{3}\) and \(23\equiv-1\pmod{4}\), while
\(29\equiv-1\pmod{3}\). Therefore \(k[R]/k\) is the unramified
quadratic extension, and \(\mathcal O[R]\) is its maximal order.

An \(R\)-stable \(\mathcal O\)-lattice in \(k^2\) is a
rank-one \(\mathcal O[R]\)-module, hence a fractional ideal of
\(\mathcal O[R]\). Since the extension is unramified, all such
lattices are homothetic under \(k^\times\). Moreover, \(\det R\) is
a unit, so a homothety class fixed by the projective transformation
\(r\) is represented by an \(R\)-stable lattice. Thus \(r\) fixes a
unique homothety class of lattices.

The subgroup \(\langle r\rangle\) is normal in \(\Gamma\), and
conjugation sends \(r\) to \(r\) or \(r^{-1}\). These elements have
the same fixed lattice classes. It follows that \(\Gamma\)
preserves the unique class above. Choose a representative
\(\Lambda\). Then
\[
 \mathcal R=\End_{\mathcal O}(\Lambda)
\]
is a \(\Gamma\)-stable maximal order in \(M_2(k)\).
Every element of \(\Gamma\) has a projective representative in
\(\GL_2(\mathcal O)\) acting on \(\Lambda\). Reduction gives a
projective action
\[
 \Gamma\longrightarrow\PGL_2(\F_p).
\]
Its natural two-dimensional projective module is irreducible over
\(\F_p\), since the reduction of \(R\) has irreducible
characteristic polynomial.

This construction uses the full endomorphism algebra \(D\).
It therefore requires no maximality assumption on
\(\End(A_{\overline{\Q}})\otimes\Z_p\).

\section{A change of Tate lattice}

Let \(V=V_p(A)\) and \(T=T_p(A)\). Define
\[
 T'=\mathcal R T.
\]
Since \(\mathcal R\) is finite free over \(\mathcal O\), this is a
full \(\mathcal O\)-lattice in \(V\). It is stable under
\(G_{\Q}=\Gal(\overline{\Q}/\Q)\), since both \(T\) and
\(\mathcal R\) are stable. After choosing a basis of \(\Lambda\),
we have \(\mathcal R=M_2(\mathcal O)\), and the matrix units give
\[
 T'\simeq\Lambda\otimes_{\mathcal O} W_0,
\]
where \(W_0=e_{11}T'\) is free of rank \(2\) over \(\mathcal O\).
The off-diagonal matrix units identify \(e_{11}T'\) and
\(e_{22}T'\), which gives the displayed decomposition.

Put \(\overline V=T'/pT'\). The reductions of two stable lattices
in the same \(p\)-adic representation have the same
semisimplification, by Brauer--Nesbitt. Since
\(T/pT=A[p]\) has a trivial submodule generated by \(P\), the module
\(\overline V\) has a trivial composition factor. Only this
composition factor is used; a fixed vector need not survive a
change of lattice.

The action of \(G_{\Q}\) on \(\overline V\) normalizes the algebra
\(M_2(\F_p)\) acting on \(U=\Lambda/p\Lambda\). Its induced
projective action on \(U\) factors through \(\Gamma\). Pull back
\[
 1\longrightarrow\F_p^\times\longrightarrow\GL_2(\F_p)
 \longrightarrow\PGL_2(\F_p)\longrightarrow1
\]
along this projective action of \(G_{\Q}\), and denote the resulting
group by \(\widetilde G\). Then \(U\) is an actual
two-dimensional \(\widetilde G\)-module. The centralizer of
\(M_2(\F_p)\) gives a two-dimensional \(\widetilde G\)-module \(W\)
such that
\[
 \overline V\simeq U\otimes_{\F_p} W.
\]
Explicitly, for \((g,B)\in\widetilde G\), the action on \(U\) is
given by \(B\), and \(B^{-1}\) times the action of \(g\) on
\(\overline V\) commutes with \(M_2(\F_p)\), so acts on the second
factor. All these representations factor through a common finite
quotient of \(\widetilde G\).

The module \(U\) is irreducible by the preceding section. The subgroup
\[
 H=\{(h,I):h\in G_L\}
\]
is normal in \(\widetilde G\) and acts trivially on \(U\). The
trivial composition factor of \(\overline V\) remains trivial on
pullback to \(\widetilde G\).

\section{The determinant contradiction}

We first show that \(W\) is irreducible. Otherwise its two composition
factors are one-dimensional. Tensoring a composition series with
\(U\), the composition factors of \(U\otimes W\) would then be
one-dimensional twists of the irreducible two-dimensional module
\(U\). There could be no trivial composition factor, a
contradiction.

Since \(H\) is normal, the restriction \(W|_H\) is semisimple.
For completeness, this form of Clifford's theorem holds in modular
characteristic: choose a simple \(H\)-submodule of \(W\). Its
\(\widetilde G\)-translates are simple \(H\)-submodules, and their
sum is a nonzero \(\widetilde G\)-submodule, hence all of \(W\).
Thus \(W|_H\) is a sum of simple \(H\)-modules.

Restriction to \(H\) preserves the existence of a trivial
composition factor. But \(\overline V|_H\) is two copies of the
semisimple module \(W|_H\), so \(W\) has a nonzero \(H\)-fixed
vector. Since \(H\) is normal, \(W^H\) is a
\(\widetilde G\)-submodule of \(W\). Irreducibility gives
\(W^H=W\). Hence \(H\) acts trivially on \(\overline V\).

Write \(\chi_p:G_{\Q}\to\F_p^\times\) for the mod-\(p\)
cyclotomic character. Every abelian variety over \(\Q\) has a
polarization over \(\Q\). The resulting nondegenerate alternating
pairing on \(V_p(A)\) shows that the determinant of its
\(p\)-adic representation is the square of the \(p\)-adic
cyclotomic character. A change of stable lattice does not change
this determinant. Reducing modulo \(p\), the trivial action of
\(H\) therefore gives
\[
 \chi_p(h)^2=1\qquad\text{for all }h\in G_L.
\]
The character \(\chi_p^2\) has order \((p-1)/2\), and its fixed field
is the maximal real cyclotomic field
\(\Q(\zeta_p+\zeta_p^{-1})\). Consequently
\[
 \Q(\zeta_p+\zeta_p^{-1})\subseteq L.
\]
For \(p=23\), this field has degree \(11\), and for \(p=29\) it has
degree \(14\). Neither degree divides \([L:\Q]\), whose only prime
factors are \(2\) and \(3\). This contradiction completes the proof
of Theorem~\ref{thm:main}.
\hfill\(\square\)

\section{Consequences for the census}

The groups
\href{https://genus-2-torsion.github.io/group.html?g=23}{\([23]\)}
and
\href{https://genus-2-torsion.github.io/group.html?g=29}{\([29]\)}
are therefore impossible in the class ``geometrically split, simple
over \(\Q\)'', as well as in the class ``split over \(\Q\)''.
The argument applies to every geometric splitting degree, including
degrees divisible by \(23\) or \(29\).

\begin{thebibliography}{9}
\bibitem{DKSS}
M. Derickx, S. Kamienny, W. Stein and M. Stoll,
\emph{Torsion points on elliptic curves over number fields of small
degree}, Algebra \& Number Theory \textbf{17} (2023), 267--308, Theorem 1.2.
\href{https://arxiv.org/pdf/1707.00364}{arXiv:1707.00364}.

\bibitem{Khawaja}
M. Khawaja,
\emph{Torsion primes for elliptic curves over degree \(8\) number fields},
Research in Number Theory \textbf{10} (2024), article 48, Theorem 1.2.
\href{https://eprints.whiterose.ac.uk/id/eprint/211212/8/s40993-024-00533-6.pdf}{Published article}.

\bibitem{Clark}
P. L. Clark,
\emph{Bounds for torsion on abelian varieties with integral moduli},
arXiv:math/0407264v2 (2004), Theorem 5.
\href{https://arxiv.org/pdf/math/0407264}{arXiv:math/0407264}.
The same list is Theorem~1.4 of P. L. Clark and X. Xarles, \emph{Local bounds for torsion points on abelian varieties}, Canad. J. Math. \textbf{60} (2008), 532--555, \href{https://doi.org/10.4153/CJM-2008-026-x}{doi:10.4153/CJM-2008-026-x}, for abelian surfaces with everywhere anisotropic reduction, which includes potentially good reduction.
\end{thebibliography}

\end{document}
